% ============================================================================
%  Loop gain (magnitude + phase, PM quoted in the legend) and PSRR vs frequency.
%  Standalone -- compile with pdflatex or tectonic.
%  Data: ac_response.csv, generated by scripts/make_ac_csv.py
%  Columns: f_hz, loop_db, loop_phase_deg, psrr_light_db, psrr_heavy_db
%
%  Conditions: ai_error_amp with back-gate-skewed event assists, 22nm FDSOI cryo models at 4 K (tt), VDD = 1.0 V, Vref = 0.8 V, C_out = 10 pF, static load so both analyses share one unambiguous operating point.
%  Loop gain is the stb analysis broken at Istb, taken at LIGHT load (~1 mA) because that is the worse case for phase margin: 61.2 deg light vs 64.6 deg heavy.
%  PSRR is measured with the ac source in series with the supply at mag=1, so |Vout| IS the transfer Vout/Vdd; the CSV negates the dB so a larger positive value means better rejection.
%
%  PHASE CONVENTION: Spectre's stb loopGain phase starts near +180 deg for negative feedback and falls, so the phase margin is simply the phase value AT the unity-gain crossing -- no 180 deg subtraction. Verified: the CSV gives 61.18 deg at 34.966 MHz and spectre's own stb.margin.stb reports PM = 61.18 deg at 34.967 MHz, agreeing to three decimals. The PM quoted in the phase legend entry is spectre's value.
%
%  Headline: loop gain 84.4 dB dc, UGF 35.0 MHz, PM 61.2 deg; PSRR 85 dB dc, 71.7 dB at 10 kHz, 51.9 dB at 100 kHz, 31.9 dB at 1 MHz, 11.5 dB at 10 MHz -- a clean 20 dB/decade rolloff.
% ============================================================================
\documentclass[border=5pt]{standalone}
\usepackage{pgfplots}
\pgfplotsset{compat=1.17}
\usepackage{siunitx}

\begin{document}
\begin{tikzpicture}

\pgfplotsset{
  width=9cm, height=5.6cm,
  xmode=log, xlabel={Frequency [\si{\hertz}]},
  grid=both, grid style={gray!25},
  tick label style={font=\small}, label style={font=\small},
  legend style={font=\small, cells={anchor=west}, draw=gray!50},
  title style={yshift=-1.2ex, font=\small},
  every axis plot/.append style={line width=0.9pt},
  xmin=1, xmax=1e9, xtick={1e0,1e2,1e4,1e6,1e8}, minor xtick={1e1,1e3,1e5,1e7,1e9},
}

% ---------------- loop gain: magnitude (left axis) + phase (right axis) ----------------
\begin{axis}[
  name=loop, xmode=log,
  title={Loop gain (light load)},
  ylabel={$|T|$ [\si{\deci\bel}]}, ymin=-60, ymax=100, ytick={-60,-40,-20,0,20,40,60,80,100},
  axis y line*=left,
  legend style={font=\small, at={(0.03,0.06)}, anchor=south west, draw=gray!50},
]
\addplot[blue] table[x=f_hz, y=loop_db, col sep=comma] {ac_response.csv};
\addlegendentry{magnitude}
\addlegendimage{red, dashed}
\addlegendentry{phase, PM $=\SI{61.2}{\degree}$}
\draw[gray!70, dotted] ({rel axis cs:0,0}|-{axis cs:1,0}) -- ({rel axis cs:1,0}|-{axis cs:1,0});   % 0 dB line
\end{axis}

\begin{axis}[
  at={(loop.south west)}, anchor=south west, xmode=log,
  width=9cm, height=5.6cm,
  ylabel={$\angle T$ [\si{\degree}]}, ymin=-60, ymax=200, ytick={-60,0,60,120,180},
  axis y line*=right, axis x line=none,
  ylabel near ticks,   % no rotate=180 here: pgfplots already turns the label for a right-hand axis, so forcing 180 renders it upside-down
]
\addplot[red, dashed] table[x=f_hz, y=loop_phase_deg, col sep=comma] {ac_response.csv};
\end{axis}

% ---------------- PSRR ----------------
\begin{axis}[
  name=psrr, at={(loop.south east)}, xshift=30mm, anchor=south west, xmode=log,   % must clear the phase axis's right-hand label, which hangs ~12mm past loop.south east
  title={Power-supply rejection},
  ylabel={PSRR [\si{\deci\bel}]}, ymin=0, ymax=100, ytick={0,20,40,60,80,100},
  legend pos=south west,
]
\addplot[blue] table[x=f_hz, y=psrr_light_db, col sep=comma] {ac_response.csv};
\addlegendentry{$\sim$\SI{1}{\milli\ampere}}
\addplot[red, dashed] table[x=f_hz, y=psrr_heavy_db, col sep=comma] {ac_response.csv};
\addlegendentry{$\sim$\SI{10}{\milli\ampere}}
\end{axis}

\end{tikzpicture}
\end{document}